\documentclass[10pt,a4paper]{amsart}
\usepackage[margin=19mm]{geometry}
\usepackage{amsmath,amssymb,mathtools}
\usepackage[scr=rsfs,cal=euler]{mathalfa}
\usepackage{xfrac}
\usepackage[unicode,hidelinks,bookmarks=false]{hyperref}
\newcommand{\ie}{i.e.\ }
\newcommand{\cf}{cf.\ }
\newcommand{\resp}{respectively\ }
\newcommand{\Subsection}{\subsection}
\providecommand{\nobreakdash}{\nobreak\hbox{-}}
\setlength{\emergencystretch}{2em}
\multlinegap0pt
\usepackage{enumitem}
\usepackage{cleveref}
\theoremstyle{plain}
\newtheorem{thm}{Theorem}
\newtheorem{lem}{Lemma}
\theoremstyle{definition}
\newtheorem{obs}{Observation}
\crefname{lem}{Lemma}{Lemmas}
\crefname{obs}{Observation}{Observations}
\newcommand{\Q}{\mathcal{Q}}
\newcommand{\G}{\mathcal{G}}
\newcommand{\C}{\mathcal{C}}
\newcommand{\FF}{\mathcal{F}}
\newcommand{\eps}{\varepsilon}
\newcommand{\sub}{\mathrm{sub}}
\newcommand{\sbb}{\mathrm{sb}}
\newcommand{\ext}{\mathrm{ext}}
\newcommand{\final}{\mathrm{final}}
\title[Finding a cycle with exactly $k$ chords]{Chi-boundedness of graphs containing no cycles with $k$ chords: finding a cycle with exactly $k$ chords from a large collection of induced $K_{\ell,\ell}$}
\author{Joonkyung Lee, Shoham Letzter, Alexey Pokrovskiy}
\date{Source: arXiv:2208.14860v3 (CC BY 4.0)}
\begin{document}
\maketitle

\noindent\emph{Source and licence.} Chi-boundedness of graphs containing no cycles with $k$ chords. Version arXiv:2208.14860v3 (CC BY 4.0), distributed under the Creative Commons Attribution 4.0 International licence, http://creativecommons.org/licenses/by/4.0/, as recorded in the arXiv licence metadata for this exact version and shown on the abstract page https://arxiv.org/abs/2208.14860v3. Full licence notice: licenses/arxiv-2208.14860-CC-BY-4.0.txt.\par\smallskip

\noindent\textit{Editorial scope.} Counted unit: the article's lemma finding a cycle with exactly $k$ chords from a large collection of pairwise vertex-disjoint induced copies of $K_{\ell,\ell}$ with no edges between them (source0.tex lines 1189--1191, source label \texttt{lem:find-cycle-from-Kll-exact}), delivered together with the article's complete original proof of it (source0.tex lines 1193--1372, one \texttt{proof} environment of 180 lines, adjacent to the statement, with no gap and no deferral). No mathematics is written, repaired or re-derived: statement and proof are the article's own text line for line, in the article's own notation ($p$-extractions, unimodal paths, layers, fathers, chords, the cycles $\C(\alpha_1,\ldots)$, the sets $X_i,Y_i$, the paths $P_{i,j},Q_i$, the cases (a)--(f)), and no retained line is substituted or re-flowed. The one adaptation is stated exactly and is bounded: the proof contains three figure environments that display the three auxiliary pictures of the proof, and the package reproduces each of those environments with its placement, its caption and its label but WITHOUT the image, so exactly three lines of the retained proof -- the three \texttt{\textbackslash includegraphics} commands, one per figure, namely the lines carrying \texttt{figures/lem23a.pdf}, \texttt{figures/lem23b.pdf} and \texttt{figures/lem23c.pdf} -- are omitted and no others; every other line of the retained spans, including the captions and labels of those three figures and every commented-out line, is carried verbatim. Required context retained uncounted: the observation that unimodal paths in different layers send no edges to one another (lines 293--296, printed as Observation 3.2); the statements (without their proofs) of the five lemmas that the counted proof invokes, namely the twenty-squares lemma (lines 660--662), the skewed-squares lemma (lines 688--690), the lemma bounding edges from one induced $K_{\ell,\ell}$ to a unimodal path (lines 763--766), its analogue for two induced copies (lines 796--800), and the cleaning-step lemma (lines 995--1004); the statement of the article's induced-matching lemma (lines 1038--1041); the statement, with its six labelled cases (a)--(f), of the article's case-analysis lemma for an induced $K_{\ell,\ell}$ (lines 1048--1074); and the subsection heading and the article's own introductory paragraph for the counted lemma (lines 1184--1187). Editorial formatting only, in addition: an AMS \texttt{amsart} preamble that restates the macros and environments the retained lines use (the macros \texttt{\textbackslash C}, \texttt{\textbackslash Q}, \texttt{\textbackslash G}, \texttt{\textbackslash FF}, \texttt{\textbackslash eps}, \texttt{\textbackslash sub}, \texttt{\textbackslash sbb}, \texttt{\textbackslash ext}, \texttt{\textbackslash final} copied from the article's own preamble, the environments \texttt{lem} and \texttt{obs} sharing one counter as in the article, the \texttt{enumitem} package for the labelled enumerations, and the \texttt{cleveref} package for the \texttt{\textbackslash Cref} commands together with the two \texttt{\textbackslash crefname} declarations that give the article's abbreviated environment names \texttt{lem} and \texttt{obs} their printed names Lemma and Observation); and numbering of the retained environments per type, so that the one retained Observation prints as Observation 1 and the eight retained Lemmas print as Lemma 1 to Lemma 8 in their order of appearance, whereas the article numbers its theorem-like environments by section; the article states its induced-matching observation, labelled \texttt{obs:induced-matching} in the source, inside a Lemma environment, and the wrapper reproduces that environment, so that statement prints as a Lemma in the wrapper as it does in the article. The article is typeset in its own \texttt{article}-class format with its own fonts and the wrapper is an amsart document. No citation occurs in any retained line, so the wrapper carries no bibliography and no bibliography entry. No figure is delivered as an image: the three figures of the counted proof are retained as captioned and labelled empty environments as described above, and the wrapper declares no asset.
    \begin{obs} \label{obs:unimodal}
        Let $P$ and $Q$ be unimodal paths in a $p$-extraction of a graph $G$ in layers $i$ and $j$, respectively, with $i < j$.
        Then the vertices in $P$ other than the end vertices and their neighbours in the path send no edges to $V(Q)$.
    \end{obs}

	\begin{lem} \label{lem:twenty-squares}
		For every $c$, every large enough $k$ can be written as a sum of exactly $20$ squares larger than $c^2$.
	\end{lem}

	\begin{lem} \label{lem:sum-skewed-squares}
		For every $c$, for every large enough $k$ which is divisible by $4$, there exist $a_1, \ldots, a_{80} \ge c$ such that $k = \sum_{i \in [80]} a_i(a_i + 1)$. 
	\end{lem}

	\begin{lem}\label{lem:sqrtk-same-Kll}
	    Let $\ell \gg k$, $p \ge 1$, and let $K$ be an induced copy of $K_{\ell, \ell}$ in a $p$-extraction of a graph~$G$.
	    If $P$ is a unimodal path starting in $K$ and $u$ is a vertex outside of $P$ with an edge into $K \setminus V(P)$, then either there is a cycle in $G$ with exactly $k$ chords or there are at most $8\sqrt k$ edges from $u$ to $P$.
	\end{lem}

    \begin{lem} \label{lem:sqrtk-different-Kll}
        Let $\ell \gg k \gg 1$, $p \ge 3$, and let
        $K$ and $K'$ be vertex-disjoint induced copies of $K_{\ell, \ell}$, with no edges between them, in a $p$-extraction of a graph $G$.
        For a unimodal path $P$ that starts in $K$ and is not in the last two layers, let $u$ be a vertex with a neighbour $w \in K' \setminus V(P)$. Then either there is a cycle in $G$ with exactly $k$ chords, or $u$ sends at most $30\sqrt{k}$ edges to $P$.
    \end{lem}

        \begin{lem} \label{lem:cleaning-step-one}
            Let $p, \ell \gg m \gg k \gg 1$ and let $G_p$ be a $p$-extraction of $G$. %with $\omega(G)<k$.
            Suppose that, in $G_p$, $K_1, \ldots, K_{\ell}$ are pairwise vertex-disjoint induced copies of $K_{\ell, \ell}$ with no edges between them. Then either there is a cycle with exactly $k$ chords, or there exist a collection $K_1', \ldots, K_m'$ of pairwise vertex-disjoint induced $K_{m, m}$'s in $G_p$ and a subset $J \subseteq [p]$ of $m$ indices such that, for every triple $T=\{j_1, j_2, j_3\} \subseteq J$ with $j_1 > j_2 > j_3$, every $i \in [m]$ and every $u \in V(K_i')$, there exist vertices $f_{1}(u;T), f_{2}(u;T), f_{3}(u;T)$ as follows:
            \begin{itemize}
                \item $f_{1}(u;T)$ is a $j_1$-father of $u$, 
                \item $f_{2}(u;T)$ is a $j_2$-father of $f_{1}(u;T)$,
                \item $f_{3}(u;T)$ is a $j_3$-father of $f_{2}(u;T)$,
                \item there are no edges between $\{f_{1}(u;T),f_{2}(u;T),f_{3}(u;T)\}$ and $\bigcup_{s \neq i}V(K_{s}')$. 
            \end{itemize}
        \end{lem}

    \begin{lem} \label{obs:induced-matching}
        Let $G_2$ be a bipartite graph on the bipartition $A\cup B$, where $|A| = |B| = m$, and let $G_1$ be its subgraph. Suppose that $d_{G_1}(a) \ge 1$ for every $a \in A$ and $d_{G_2}(b) \le d$ for every $b \in B$.
        Then there exist $A' \subseteq A$ and $B' \subseteq B$, such that $|A'| = |B'| \ge m/4d^2$ and each induced subgraph $G_i[A', B']$, $i=1,2$, is a perfect matching. 
    \end{lem}%\vspace{-5mm}

    \begin{lem}\label{lem:cases}
        For $\ell \gg \omega$, let $K$ be an induced copy of $K_{\ell,\ell}$ with bipartition $X\cup Y$.
        Suppose that $R, S, T$ are pairwise disjoint sets outside $K$ such that every vertex in $K$ has a neighbour in~$R$, every vertex in $R$ has a neighbour in $S$, and every vertex in $S$ has a neighbour in $T$. 
        Let $U:=R\cup S\cup T$ and $\ell' := \log \ell / 10000$.
        
        
        Then there exist sets $X'$ and $Y'$ of size $\ell'$, each of which contained in a different set amongst $X$ and $Y$, sets $A, A', A'' \subseteq U$, each of which contained in a different set amongst $R, S, T$, and vertices $a, b \in U$ such that one of the following holds:
        % do we care about A, A', A'' coming from different sets R, S, T?
        \begin{enumerate} [label = (K\arabic*)]
            \item \label{itm:a}
                The vertex $a$ is complete to $X' \cup Y'$.
            \item \label{itm:b}
                The vertex $a$ is complete to $X'$ and anti-complete to $Y'$, and $b$ is complete to $Y'$ and anti-complete to $X'$. 
            \item \label{itm:c}
                Both the induced bipartite graphs $G[A, X']$ and $G[A', Y']$ are perfect matchings, $A$ is anti-complete to $Y'$, and $A'$ is anti-complete to $X'$.
                
            \item \label{itm:d}
                The induced bipartite graph $G[A, X']$ is a perfect matching and $A$ is complete to $Y'$.
            \item \label{itm:e}
                The induced bipartite graph $G[A, X']$ is a perfect matching, $A$ is anti-complete to $Y'$, $A'$ is anti-complete to $X'$ and either complete or anti-complete to $Y'$, every vertex in $A$ has a neighbour in $A'$, and $b$ is complete to $Y'$ and anti-complete to $X'$.
                
            \item \label{itm:f}
                Each component in $G[A, A', A'', X']$ consists of four vertices $u, u', u'', x$ in $A, A', A'', X'$, respectively, such that all pairs in $\{u, u', u'', x\}$, except for possibly $uu''$, form edges. Additionally, $A \cup A' \cup A''$ is anti-complete to $Y'$, and $b$ is complete to $Y'$ and anti-complete to $X'$.
                
                %need to change below
        \end{enumerate}
    \end{lem}

    \subsection{Finding a cycle with exactly $k$ chords}
    
    Finally, we prove the following lemma, which finds a cycle with exactly $k$ chords given a large collection of pairwise vertex-disjoint $K_{\ell,\ell}$'s with no edges between them.
    As mentioned previously, we will use a similar approach to that used in the previous section and earlier in this section, connecting the $K_{\ell,\ell}$'s by unimodal paths and closing a cycle by choosing a path of the right length from each $K_{\ell,\ell}$. The difference here is the much more careful analysis of the interaction between a single $K_{\ell,\ell}$, and a collection of parents, `grandparents' and `great-grandparents' of its vertices. This analysis will give rise to three cases (the first corresponding to the first four cases in \Cref{lem:cases} and the last two each corresponding to one of the last two cases in \Cref{lem:cases}). In the last two cases we will sometimes need to adjust the cycle lengths slightly. For that we build our cycle so as to contain certain small `gadgets' that will allows us to do so.  

    \begin{lem} \label{lem:find-cycle-from-Kll-exact}
        Let $p, \ell \gg k \gg 1$. Suppose that, in a $p$-extraction of $G$, there are $\ell$ pairwise vertex-disjoint induced copies of $K_{\ell, \ell}$ with no edges between them. Then there is a cycle in $G$ with exactly~$k$ chords.
    \end{lem}

    \begin{proof}
        Suppose that there is no cycle with exactly $k$ chords. 
        Let $m$ satisfy $p, \ell \gg m \gg k$. 
        Then by \Cref{lem:cleaning-step-one}, there exist a collection $K_1, \ldots, K_{200}$ of pairwise vertex-disjoint copies of induced $K_{m,m}$'s in the $p$-extraction of $G$, a subset $J \subseteq [p]$ of size $200 \cdot 75$ that satisfy:
        for every triple $T=\{j_1,j_2,j_3\}$ in $J$ with $j_1>j_2>j_3$, every $i\in[200]$ and every $u\in K_i$,
        there exist vertices $f_{1}(u;T)$, $f_{2}(u;T)$,  and $f_{3}(u;T)$
        such that $f_t(u;T)$ is $j_t$-father of $f_{t-1}(u;T)$, $t=1,2,3$, where $u=f_0(u;T)$. Moreover, there are no edges between $\{f_{1}(u;T),f_{2}(u;T),f_{3}(u;T)\}$ and $\cup_{s\neq i}V(K_s)$.
        Let $S_1, \ldots, S_{200}\subseteq J$ be pairwise disjoint subsets of size $75$.

        Fix $i \in [200]$ and let $T_{1}, \ldots, T_{25}$ be disjoint triples that partition $S_i$. Now repeatedly apply \Cref{lem:cases} twenty five times to each $K_i$: at the beginning, let $K=K_i$ with the bipartition $X_i^{(0)}\cup Y_i^{(0)}$. At the $j$-th iteration, we apply \Cref{lem:cases} with $X=X_i^{(j-1)}$, $Y=Y_i^{(j-1)}$, $R=\{f_{1}(u;T_j) : u \in X_i^{(j-1)} \cup Y_i^{(j-1)}\}$, $S=\{f_{2}(u;T_j) : u \in X_i^{(j-1)} \cup Y_i^{(j-1)}\}$ and $T=\{f_{3}(u;T_j) : u \in X_i^{(j-1)} \cup Y_i^{(j-1)}\}$.
        Denote the output of the lemma by $X'_{j}, Y'_{j}, A_{j}, A'_{j}, A''_{j}$ and vertices $a_{j},b_{j}$. Recall that either $X'_{j} \subseteq X$ and $Y'_{j} \subseteq Y$, or $X'_{j} \subseteq Y$ and $Y'_{j} \subseteq X$.
        In the former case, take $X_i^{(j)} = X_{j}'$ and $Y_i^{(j)} = Y_{j}'$ and otherwise take $X_i^{(j)} = Y_{j}'$ and $Y_i^{(j)} = X_{j}'$; so that $X_i^{(j)} \subseteq X_{i-1}^{(j)}$ and $Y_i^{(j)} \subseteq Y_{i-1}^{(j)}$. 

        Notice that there exists $L \subseteq [25]$ of size $3$ such that the same case among \ref{itm:a} to \ref{itm:f} occurred in all iterations indexed by $L$, and moreover either for all $\ell \in L$ we have $X_i^{(\ell)} = X_{\ell}'$ or we always have $X_i^{(\ell)} = Y_{\ell}'$.
		We assume that $X_i^{(\ell)} = X_\ell'$ for all $\ell \in L$; the other case can be handled analogously.
		We consider three cases, as follows.
        For brevity, write $X_i = X_i^{(25)}$ and $Y_i = Y_i^{(25)}$; so $X_i$ and $Y_i$ are the shrunken sets at the end of the 25 iterations, each of which has size at least $(\log ^{(50)} m) / 10^6$.\footnote{Here 
        $\log ^{(\ell)} x$ denotes $\ell$ compositions of logarithm, e.g., $\log^{(2)}x = \log\log x$.}
        %By \Cref{lem:cases} and by possibly shrinking $K_i$ (such that its parts have size at least $(\log ^{(26)} m) / 1000$),\footnote{Here 
        %$\log ^{(\ell)} x$ denotes $\ell$ compositions of logarithm, e.g., $\log^{(2)}x = \log\log x$.}
        %one of the following holds, where $\{X_i, Y_i\}$ is the bipartition of $K_i$ (after shrinking). 
        \begin{enumerate} [label = \rm(\arabic*)]
            \item \label{itm:easy-case}
				One of \ref{itm:a} to \ref{itm:d} occurred in all iterations indexed by $L$.

                Let $\ell_1, \ell_2 \in L$ be distinct.
                We choose vertices $u_{i,1}, u_{i,2}, u_{i,1}', u_{i,2}'$ as follows.
                \begin{itemize}
                    \item \ref{itm:a}:
                        Take $u_{i, 1} = a_{\ell_1}$ and $u_{i, 2} = a_{\ell_2}$, so that $u_{i,1}$ and $u_{i,2}$ are complete to $X_i \cup Y_i$; choose $u_{i, 1}' \in X_i$ and $u_{i, 2}' \in Y_i$. 
                    \item \ref{itm:b}:
						Take $u_{i,1} = a_{\ell_1}$ and $u_{i,2} = b_{\ell_2}$, so that $u_{i,1}$ is complete to $X_i$ and anti-complete to $Y_i$ and $u_{i, 2}$ is complete to $Y_i$ and anti-complete to $X_i$; choose $u_{i, 1}' \in X_i$ and $u_{i, 2}' \in Y_i$.
                    \item \ref{itm:c}:
						Pick $u_{i,1}' \in X_i$ and $u_{i,2}' \in Y_i$ arbitrarily, and let $u_{i,1}$ be the unique neighbour of $u_{i,1}'$ in $A_{\ell_1}$ and $u_{i,2}$ to be the unique neighbour of $u_{i,2}'$ in $A'_{\ell_2}$.
						Notice that $u_{i,s}'$ is the unique neighbour of $u_{i,s}$ in $X_i \cup Y_i$, for $s \in [2]$.
                    \item \ref{itm:d}:
						Let $u_{i,1}', u_{i,2}' \in X_i$ be distinct, and let $u_{i,s}$ be unique neighbour of $u_{i,s}'$ in $A_{\ell_s}$, for $s \in [2]$.
						Then $u_{i,s}$ is complete to $Y_i$ and $u_{i,s}'$ is its only neighbour in $X_i$, for $s \in [2]$.
                \end{itemize}
				Notice that $u_{i,1},u_{i,2}$ belong to different layers with indices in $S_i$ and are anti-complete to $X_j$ for $j \in [200] \setminus \{i\}$.
            \item \label{itm:annoying-case-a}
				Case \ref{itm:e} occurred in all cases indexed by $L$.

				Write $L = \{\ell_1, \ell_2, \ell_3\}$.
				Let $u_{i,1}', u_{i,3}' \in X_i$ be distinct and let $u_{i,s}$ be the unique neighbour of $u_{i,s}'$ in $A_{\ell_s}$, for $s \in \{1,3\}$. Take $u_{i,4}$ to be a neighbour of $u_{i,3}$ in $A'_{\ell_3}$, take $u_{i,2} = b_{\ell_2}$ and choose $u_{i,2}' \in Y_i$.
				Then
                \begin{itemize}
                    \item
						$u_{i, 1}, \ldots, u_{i, 4}$ belong to four distinct layers with indices in $S_i$ and are anti-complete to $X_j \cup Y_j$, for $j \in [200] \setminus \{i\}$.
                    \item 
						$u_{i, 1}$ and $u_{i, 3}$ are anti-complete to $Y_i$, and $u_{i,s}'$ is the unique neighbour of $u_{i,s}$ in $X_i$, for $s \in \{1,3\}$.
                    \item
                        $u_{i, 2}$ is complete to $Y_i$ and anti-complete to $X_i$. 
                    \item
                        $u_{i, 4}$ is anti-complete to $X_i$ and either complete or anti-complete to $Y_i$.
                \end{itemize}
            \item \label{itm:annoying-case-b}
				Case \ref{itm:f} occurred in all cases indexed by $L$.

				Write $L = \{\ell_1, \ell_2, \ell_3\}$.
				Let $u_{i,1}', u_{i,3}' \in X_i$ be distinct and let $u_{i,s}$ be the unique neighbour of $u_{i,s}'$ in $A_{\ell_s}$, for $s \in \{1,3\}$. Take $u_{i,4}$ to be the unique neighbour of $u_{i,3}$ in $A'_{\ell_3}$ and $u_{i,5}$ the unique neighbour of $u_{i,4}$ in $A''_{\ell_3}$. Finally, let $u_{i,2} = b_{\ell_2}$ and pick $u_{i,2}' \in Y_i$.
				Then
                \begin{itemize}
                    \item
						$u_{i, 1}, \ldots, u_{i, 5}$ are in distinct layers with indices in $S_i$ and are anti-complete to $X_j \cup Y_j$ for $j \in [200] \setminus \{i\}$.
                    \item 
						$u_{i,s}'$ is the unique neighbour of $u_{i,s}$ in $X_i$ and is anti-complete to $Y_i$, for $s \in \{1,3\}$.
                    \item
                        $u_{i, 2}$ is complete to $Y_i$ and anti-complete to $X_i$. 
                    \item
                        $u_{i, 4}$ is a neighbour of $u_{i, 3}$ and $u_{i, 3}'$ and is anti-complete to $(X_i \cup Y_i) \setminus \{u_{i, 3}'\}$.
                    \item
                        $u_{i, 5}$ is a neighbour of $u_{i, 3}'$ and $u_{i, 4}$ and is anti-complete to $(X_i \cup Y_i) \setminus \{u_{i, 3}'\}$.
                        % now u_{i,3} and u_{i,5} need not be adjacent
                \end{itemize}
        \end{enumerate}
        %It follows that at least one of \ref{itm:easy-case}, \ref{itm:annoying-case-a} and \ref{itm:annoying-case-b} holds.
        Notice that one of the following holds: \ref{itm:easy-case} holds for $20$ values of $i \in [200]$; \ref{itm:annoying-case-a} holds for $100$ values of $i \in [200]$; or \ref{itm:annoying-case-b} holds for $80$ values of $i \in [200]$.
		We show how to find the desired cycle with $k$ chords in each of these cases separately.
        
        \medskip
        
        \textbf{Case \ref{itm:easy-case} holds for at least $20$ values of $i \in [200]$.}
        
        By relabelling the indices of the $200$ copies of $K_{m,m}$ if necessary, we may assume that \ref{itm:easy-case} holds for $i \in [20]$. To construct a cycle with $k$ chords, we need one more copy of $K_{m,m}$, so we shall use $K_{21}$ too. For simplicity, let $K_0=K_{21}$ be the copy of $K_{m,m}$ on $X_0\cup Y_0$, where $X_0=X_{21}$ and $Y_0=Y_{21}$.
        
		We want to take a set of distinct vertices $w_{i,j}$, $i\in[20]$ and $j \in [2]$, where each $w_{i,j}$ is at the same layer as $u_{i,j}$ and has a neighbour $w_{i,j}'$ in $K_0$ which are all distinct too, and there are no edges between $\bigcup_{i \in [20]} (X_i \cup Y_i)$ and $\{w_{i,j}: i \in [20], j \in [2]\}$. Here the index $j=1,2$ indicates $w_{i,1}'\in X_0$ and $w_{i,2}'\in Y_0$, as was done for $u_{i,j}'$'s.
        Indeed, this is possible since one can greedily take $w_{i,1}'$ and $w_{i,2}'$ disjoint from the previous choices in $X_0$ and $Y_0$, respectively, and let $w_{i,j}=f_1(w'_{i,j};T)$ where $T$ is any triple whose smallest index is the index of the layer containing $u_{i,j}$.
        %let $w_{i, j}$ be a vertex at the same layer as $u_{i, j}$ such that each of $w_{i, 1}$ and $w_{i,2}$ has a neighbour $w_{i, 1}'$ in $X_{0}$ and $w_{i,2}'$ in $Y_{0}$, respectively, and there are no edges between $\{w_{i,1},w_{i,2}\}$ and $\cup_{i=1}^{20}V(K_i)$. Such vertices $w_{i,j}$ and $w_{i,j}'$ exist as we can take a vertex in $K_{0}$. 
        %We further require that the vertices $w_{i, j}'$, with $j \in [2]$ and $i \in [20]$, are all distinct. Such vertices $w_{i, j}, w_{i, j}'$ exist by the assumptions on $K_0, \ldots, K_{20}$. 
        Let $P_{i, j}$ be a unimodal path between $u_{i, j}$ and $w_{i, j}$, taken in the same layers as the two end vertices so that they are in different layers from the $K_i$'s. 

		Given integers $\alpha_1, \ldots, \alpha_{20}$, we define a cycle $\C(\alpha_1, \ldots, \alpha_{20})$ as follows. Let $Q_i$ be the path from $w_{i, 1}'$ to $w_{i+1, 1}$ for $i \in [20]$ (addition of indices taken modulo $20$), obtained by concatenating the paths $w_{i, 1}' w_{i, 1}$, $P_{i,1}$, $u_{i, 1} u_{i, 1}'$, $Q_i^*$, $u_{i, 2}' u_{i, 2}$, $P_{i,2}$, $w_{i, 2} w_{i, 2}' w_{i+1, 1}'$, where $Q_i^*$ is a path in $G[X_i, Y_i]$ with ends $u_{i, 1}'$ and $u_{i, 2}'$ of length $2\alpha_i + 1$. 
        Now let $\C(\alpha_1, \ldots, \alpha_{20})$ be the cycle obtained by concatenating $Q_1, \ldots, Q_{20}$.
        
        \begin{figure}[h]
            \centering

            \caption{The path $Q_i$ in Case \ref{itm:easy-case}}
            \label{fig:easy-case}
        \end{figure}
        
        While the cycle $\C(\alpha_1, \ldots, \alpha_{20})$ is not uniquely defined, its number of chords depends only on $\alpha_1, \ldots, \alpha_{20}$. 
        To evaluate the number of chords in this cycle, $Q_i^*$ contributes $(\alpha_i + 1)^2 - (2\alpha_i + 1) = \alpha_i^2$ chords; the number of chords with one end in $\{u_{i, 1}, u_{i, 2}\}$ and one end in $X_i \cup Y_i$ is one of the four values $2(2\alpha_i + 1)$, $2\alpha_i$, $0$, and $2(\alpha_i + 1)$, depending on the four cases described in~\ref{itm:easy-case}; there are at most $120^2$ chords with both ends in $\bigcup_{i \in [20], j \in [2]} \{w_{i, j}, w_{i, j}', u_{i, j}\}$; and we may assume that the number of chords with ends in $\bigcup_{i \in [20], j \in [2]} V(P_{i,j})$ is $O(\sqrt{k})$ by \Cref{lem:sqrtk-different-Kll,lem:sqrtk-same-Kll}. Note that there are no chords between distinct sets $X_i \cup Y_i$, between a set $X_i \cup Y_i$ and $\{u_{j,1}, u_{j,2}, w_{j, 1}, w_{j, 2}\}$ for distinct $i, j \in [20]$, or between $\{u_{i,1}, u_{i, 2}\}$ and $\{w_{j, 1}, w_{j,1}', w_{j,2}, w_{j,2}'\}$.  In total, the number of chords is 
        \begin{equation*}
            f_G(k) + \sum_{i \in [20]} (\alpha_i^2 + 2\sigma_i \alpha_i) 
            = f_G(k) + \sum_{i \in [20]} \left((\alpha_i + \sigma_i)^2 - \sigma_i^2\right) 
            = h_G(k) + \sum_{i \in [20]} (\alpha_i + \sigma_i)^2.
        \end{equation*}
        where $\sigma_i \in \{0, 1, 2\}$ for $i \in [20]$, and $f_G(k)$ and $h_G(k)$ are functions that do not depend on $\alpha_1, \ldots, \alpha_{20}$ and satisfy $f_G(k), h_G(k) = O(\sqrt{k})$. By \Cref{lem:twenty-squares}, there is a choice of integers $\beta_1 \ldots, \beta_{20} \ge 2$ such that $\sum_{i \in [20]} \beta_i^2 = k - h_G(k)$. A cycle $\C(\alpha_1, \ldots, \alpha_{20})$ with $\alpha_i = \beta_i - \sigma_i$ has $k$ chords, as desired.
        
        \medskip
        
        \textbf{Case \ref{itm:annoying-case-a} holds for at least $100$ values of $i \in [200]$.}
        
        Again by possible relabelling, we may assume that \ref{itm:annoying-case-a} holds for $i \in [100]$.
        We also need an extra copy, denoted by $K_0$, taken from $K_i$, $i>100$, on the bipartition $X_0\cup Y_0$.
        Let $w_{i, j}$, with $i \in [100]$ and $j \in [4]$, satisfy the following: $w_{i, j}$ is at the same layer as $u_{i, j}$; it is anti-complete to $X_i \cup Y_i$ for $i \in [100]$; $w_{i, j}$ has a neighbours $w_{i, j}'$ such that $w_{i, j}' \in X_0$ if $j \in \{1, 3\}$ and otherwise $w_{i, j}' \in Y_0$; and the vertices $w_{i, j}, w_{i, j}'$, with $j \in [4]$ and $i \in [100]$, are all distinct. 
        Indeed, such choices of $w_{i,j}$ and $w_{i,j}'$ are possible since one can again take $w_{i,j}=f_1(w_{i,j}';T)$ with any triple $T$ whose smallest index is the index of the layer containing $u_{i,j}$, while maintaining all $w_{i,j}'$ being distinct.
        
        Let $P_{i, j}$ be a unimodal path with ends $u_{i, j}$ and $w_{i, j}$ in the same layer as the end vertices,
        for each $i \in [100]$ and $j \in [4]$. 
        Given positive integers $\alpha_1, \ldots, \alpha_{100}$, we define a cycle $\C(\alpha_1, \ldots, \alpha_{100})$ as follows. Let $Q_i$ be the path from $w_{i, 1}'$ to $w_{i+1, 1}'$ (index addition taken modulo $100$), obtained by concatenating the paths $w_{i, 1}' w_{i, 1}$, $P_{i, 1}$, $u_{i, 1} u_{i, 1}'$, $Q_i^*$, $u_{i, 2}'u_{i, 2}$, $P_{i, 2}$, $w_{i, 2} w_{i, 2}' w_{i, 3}' w_{i, 3}$, $P_{i,3}$, $u_{i, 3} u_{i, 4}$, $P_{i, 4}$, $w_{i, 4} w_{i, 4}' w_{i+1, 1}'$, where $Q_{i}^*$ is a path in $G[X_i \setminus \{u_{i, 3}'\}, Y_i]$ with ends $u_{i, 1}'$ and $u_{i, 2}'$, and of length $2\alpha_i + 1$. Let $\C(\alpha_1, \ldots, \alpha_{100})$ be the cycle obtained by concatenating $Q_1, \ldots, Q_{100}$.
        
        \begin{figure}
            \centering

            \caption{The path $Q_i$ in Case \ref{itm:annoying-case-a}}
            \label{fig:annoying-case-a}
        \end{figure}
        
        As above, the cycle $\C(\alpha_1, \ldots, \alpha_{100})$ is not uniquely defined, but the number of its chords depends only on $\alpha_1, \ldots, \alpha_{100}$. 
        Let us evaluate the number of chords precisely as follows: 
        $Q_{i}^*$ contributes $\alpha_i^2$ chords; 
        the number of chords between $\{u_{i, 1}, \ldots, u_{i, 4}\}$ and $X_i \cup Y_i$ is either $\alpha_i$ or $2\alpha_i+1$ depending on how $u_{i,4}$ connects to $Y_i$ (either complete or anti-complete); there are at most $1200^2$ chords with ends in $\bigcup_{i \in [100], j \in [4]}\{w_{i, j}, w_{i, j}', u_{i, j}\}$; and we may assume that the number of chords in $\bigcup_{i \in [100], j \in [4]} V(P_{i, j})$ is $O(\sqrt{k})$ by \Cref{lem:sqrtk-different-Kll,lem:sqrtk-same-Kll}. In total, the number of chords in the cycle is
        \begin{align*}
            f_G(k) + \sum_{i \in [100]}(\alpha_i^2 + \sigma_i \alpha_i),
        \end{align*}
        where $\sigma_i \in \{1, 2\}$ and $f_G(k) = O(\sqrt{k})$. Here $\sigma_i$'s and $f_G(k)$ do not depend on $\alpha_i$'s.
        
        Suppose that $\sigma_i = 2$ for at least $20$ values of $i \in [100]$; let $I$ be a set of $20$ such values of $i$. Set $\alpha_i = 1$ for $i \in [100] \setminus I$. Each $\alpha_i$, $i \in I$, is not yet determined. Then the number of chords in a cycle $\C(\alpha_1, \ldots, \alpha_{100})$ reduces to 
        \begin{equation*}
            h_G(k) + \sum_{i \in I} (\alpha_i + 1)^2,
        \end{equation*}
        for some $h_G(k) = O(\sqrt{k})$ that does not depend on $\alpha_i$'s with $i \in I$. By \Cref{lem:twenty-squares}, there exist integers $\beta_i$, for $i \in I$, such that $\beta_i \ge 2$ and $\sum_{i \in I} \beta_i^2 = k - h_G(k)$. Set $\alpha_i = \beta_i - 1$. The cycle $\C(\alpha_1, \ldots, \alpha_{100})$ then has exactly $k$ chords.
        
        It remains to consider the case where $\sigma_i = 1$ for at least $80$ values of $i \in [100]$. Let $I$ be a set of $80$ values of $i$ such that $\sigma_i=1$. Set $\alpha_i = 1$ for $i \in [100] \setminus I$. Then the number of chords in the cycle $\C(\alpha_1, \ldots, \alpha_{100})$ is 
        \begin{equation*} 
             h_G(k) + \sum_{i \in I} (\alpha_i^2 + \alpha_i),
        \end{equation*}
        where $h_G(k) = O(\sqrt{k})$ is independent of the $\alpha_i$'s. Let $\tau \in \{0, 1, 2, 3\}$ be the remainder of $k - h_G(k)$ modulo $4$. Let $\C_{\tau}(\alpha_1, \ldots, \alpha_{100})$ be a cycle obtained by replacing an inner vertex of $V(Q_i^*) \cap X_i$ by $u_{i, 3}'$, for $\tau$ values of $i$. One can readily check that the number of chords in $\C_{\tau}(\alpha_1, \ldots, \alpha_{100})$ is $h_G(k) + \tau + \sum_{i \in I} (\alpha_i^2 + \alpha_i)$m as we gain the additional chord $u_{i, 3} u_{i, 3}'$ for $\tau$ values of $i$. Now, since $k - h_G(k) - \tau$ is divisible by $4$, \Cref{lem:sum-skewed-squares} implies that there exist integers $\alpha_i$, $i \in I$, such that $\alpha_i \ge 1$ and $\sum_{i \in I} (\alpha_i^2 + \alpha_i) = k - h_G(k) - \tau$. In particular, there is a cycle with exactly $k$ chords.
        
        \medskip
        
        \textbf{Case \ref{itm:annoying-case-b} holds for at least $80$ values of $i \in [200]$}
        
        Without loss of generality, \ref{itm:annoying-case-b} holds for $i \in [80]$ and let $K_0=K_{81}$ be another copy of $K_{m,m}$ on $X_0\cup Y_0$. Analogously to the previous cases, choose $w_{i, j}$, $i \in [80]$ and $j \in [5]$, that satisfy the following conditions: $w_{i, j}$ is in the same layer as $u_{i, j}$; it is anti-complete to $X_i \cup Y_i$ for $i \in [80]$; $w_{i, j}$ has a neighbour $w_{i, j}'$ such that $w_{i, j}' \in X_0$ for $j \in \{1, 3\}$ and $w_{i, j}' \in Y_0$ for $j \in \{2, 4, 5\}$; and the vertices $w_{i, j}, w_{i,j}'$, with $i \in [80]$ and $j \in [5]$, are all distinct. 
        As before, such choices for $w_{i,j}$ and $w_{i,j}'$ exist.
        
        One can find a path $P_{i, j}$ between $u_{i, j}$ and $w_{i, j}$ in the same layer as the two ends, for each $i \in [80]$ and $j \in [5]$. 
        In fact, we do not make use of $w_{i, 4}$, $w_{i, 4}'$ and $P_{i, 4}$, but we took them to preserve the correspondence between indices.
        
        Given positive integers $\alpha_1, \ldots, \alpha_{80}$, we choose a cycle $\C^0(\alpha_1, \ldots, \alpha_{80})$ as follows. Let $Q_i$ be the path from $w_{i, 1}'$ to $w_{i+1, 1}'$ (addition modulo $80$), obtained by concatenating the paths $w_{i,1}' w_{i,1}$, $P_{i, 1}$, $u_{i, 1} u_{i, 1}'$, $Q_{i}^*$, $u_{i, 2}' u_{i, 2}$, $P_{i, 2}$, $w_{i, 2} w_{i, 2}' w_{i, 3}' w_{i, 3}$, $P_{i, 3}$, $u_{i, 3} u_{i, 4} u_{i, 5}$, $P_{i, 5}$, $w_{i, 5} w_{i, 5}' w_{i+1, 1}'$; where $Q_i^*$ is a path in $G[X_i \setminus \{u_{i, 3'}\}, Y_i]$ with ends $u_{i,1}'$ and $u_{i, 2}'$ and of length $2\alpha_i + 1$. Let $\C^0(\alpha_1, \ldots, \alpha_{80})$ be the cycle obtained by concatenating $Q_1, \ldots, Q_{80}$. For $\tau =1,2,3$, let $\C^{\tau}(\alpha_1, \ldots, \alpha_{80})$ be the cycle obtained from $\C^0(\alpha_1, \ldots, \alpha_{80})$ by replacing one of the internal vertices of $Q_{i}^*$ in $X_i$ by $u_{i, 3}'$, for $i \in [\tau]$. 
        \begin{figure}
            \centering

            \caption{The path $Q_i$ in Case \ref{itm:annoying-case-b}}
            \label{fig:annoying-case-b}
        \end{figure}
        
        The cycle $\C^\sigma(\alpha_1, \ldots, \alpha_{80})$ for each $\sigma=0,1,2,3$ is not necessarily unique, but all the choices have the same number of chords. 
        To verify, let us evaluate the number of chords in $\C^0(\alpha_1, \ldots, \alpha_{80})$ first, as follows: $Q_i^*$ contributes $\alpha_i^2$ chords; there are exactly $\alpha_i$ chords between $\{u_{i,1}, \ldots, u_{i, 5}\}$ and $X_i \cup Y_i$, all of which are incident to $u_{i,2}$; there are at most $1200^2$ chords with ends in $\bigcup_{i \in [80], j \in [5]} \{w_{i, j}, w_{i,j}', u_{i, j}\}$; and we may assume that the number of chords in $\bigcup_{i \in [80], j \in [5]} V(P_{i, j})$ is $O(\sqrt{k})$, by \Cref{lem:sqrtk-different-Kll,lem:sqrtk-same-Kll}. 
        The cycle $\C^{\tau}(\alpha_1, \ldots, \alpha_{80})$ has precisely $3\tau$ more chords than $\C^0(\alpha_1, \ldots, \alpha_{80})$, as $u_{i,3}'$ gives three extra chords to $u_{i,3}$, $u_{i,4}$, and $u_{i,5}$ (see \Cref{fig:annoying-case-b}).
        The total number of chords in $\C^{\tau}(\alpha_1, \ldots, \alpha_{80})$ is hence
        \begin{equation*}
            f_G(k) + \sum_{i \in [80]}(\alpha_i^2 + \alpha_i) + 3\tau,
        \end{equation*}
        where $f_G(k) = O(\sqrt{k})$. 
        Now choose $\tau \in \{0, 1, 2, 3\}$ be such that $3\tau$ equals $k - f_G(k)$ modulo $4$. By \Cref{lem:sum-skewed-squares}, there exist positive integers $\alpha_1, \ldots, \alpha_{80}$, such that $\sum_{i \in [80]} \alpha_i(\alpha_i + 1) = k - f_G(k) - 3\tau$. The cycle $\C^{\tau}(\alpha_1, \ldots, \alpha_{80})$ has exactly $k$ chords.
    \end{proof}

\end{document}
